% Options for packages loaded elsewhere
% Options for packages loaded elsewhere
\PassOptionsToPackage{unicode}{hyperref}
\PassOptionsToPackage{hyphens}{url}
\PassOptionsToPackage{dvipsnames,svgnames,x11names}{xcolor}
%
\documentclass[
  letterpaper,
  DIV=11,
  numbers=noendperiod]{scrartcl}
\usepackage{xcolor}
\usepackage{amsmath,amssymb}
\setcounter{secnumdepth}{-\maxdimen} % remove section numbering
\usepackage{iftex}
\ifPDFTeX
  \usepackage[T1]{fontenc}
  \usepackage[utf8]{inputenc}
  \usepackage{textcomp} % provide euro and other symbols
\else % if luatex or xetex
  \usepackage{unicode-math} % this also loads fontspec
  \defaultfontfeatures{Scale=MatchLowercase}
  \defaultfontfeatures[\rmfamily]{Ligatures=TeX,Scale=1}
\fi
\usepackage{lmodern}
\ifPDFTeX\else
  % xetex/luatex font selection
\fi
% Use upquote if available, for straight quotes in verbatim environments
\IfFileExists{upquote.sty}{\usepackage{upquote}}{}
\IfFileExists{microtype.sty}{% use microtype if available
  \usepackage[]{microtype}
  \UseMicrotypeSet[protrusion]{basicmath} % disable protrusion for tt fonts
}{}
\makeatletter
\@ifundefined{KOMAClassName}{% if non-KOMA class
  \IfFileExists{parskip.sty}{%
    \usepackage{parskip}
  }{% else
    \setlength{\parindent}{0pt}
    \setlength{\parskip}{6pt plus 2pt minus 1pt}}
}{% if KOMA class
  \KOMAoptions{parskip=half}}
\makeatother
% Make \paragraph and \subparagraph free-standing
\makeatletter
\ifx\paragraph\undefined\else
  \let\oldparagraph\paragraph
  \renewcommand{\paragraph}{
    \@ifstar
      \xxxParagraphStar
      \xxxParagraphNoStar
  }
  \newcommand{\xxxParagraphStar}[1]{\oldparagraph*{#1}\mbox{}}
  \newcommand{\xxxParagraphNoStar}[1]{\oldparagraph{#1}\mbox{}}
\fi
\ifx\subparagraph\undefined\else
  \let\oldsubparagraph\subparagraph
  \renewcommand{\subparagraph}{
    \@ifstar
      \xxxSubParagraphStar
      \xxxSubParagraphNoStar
  }
  \newcommand{\xxxSubParagraphStar}[1]{\oldsubparagraph*{#1}\mbox{}}
  \newcommand{\xxxSubParagraphNoStar}[1]{\oldsubparagraph{#1}\mbox{}}
\fi
\makeatother


\usepackage{longtable,booktabs,array}
\newcounter{none} % for unnumbered tables
\usepackage{calc} % for calculating minipage widths
% Correct order of tables after \paragraph or \subparagraph
\usepackage{etoolbox}
\makeatletter
\patchcmd\longtable{\par}{\if@noskipsec\mbox{}\fi\par}{}{}
\makeatother
% Allow footnotes in longtable head/foot
\IfFileExists{footnotehyper.sty}{\usepackage{footnotehyper}}{\usepackage{footnote}}
\makesavenoteenv{longtable}
\usepackage{graphicx}
\makeatletter
\newsavebox\pandoc@box
\newcommand*\pandocbounded[1]{% scales image to fit in text height/width
  \sbox\pandoc@box{#1}%
  \Gscale@div\@tempa{\textheight}{\dimexpr\ht\pandoc@box+\dp\pandoc@box\relax}%
  \Gscale@div\@tempb{\linewidth}{\wd\pandoc@box}%
  \ifdim\@tempb\p@<\@tempa\p@\let\@tempa\@tempb\fi% select the smaller of both
  \ifdim\@tempa\p@<\p@\scalebox{\@tempa}{\usebox\pandoc@box}%
  \else\usebox{\pandoc@box}%
  \fi%
}
% Set default figure placement to htbp
\def\fps@figure{htbp}
\makeatother


% definitions for citeproc citations
\NewDocumentCommand\citeproctext{}{}
\NewDocumentCommand\citeproc{mm}{%
  \begingroup\def\citeproctext{#2}\cite{#1}\endgroup}
\makeatletter
 % allow citations to break across lines
 \let\@cite@ofmt\@firstofone
 % avoid brackets around text for \cite:
 \def\@biblabel#1{}
 \def\@cite#1#2{{#1\if@tempswa , #2\fi}}
\makeatother
\newlength{\cslhangindent}
\setlength{\cslhangindent}{1.5em}
\newlength{\csllabelwidth}
\setlength{\csllabelwidth}{3em}
\newenvironment{CSLReferences}[2] % #1 hanging-indent, #2 entry-spacing
 {\begin{list}{}{%
  \setlength{\itemindent}{0pt}
  \setlength{\leftmargin}{0pt}
  \setlength{\parsep}{0pt}
  % turn on hanging indent if param 1 is 1
  \ifodd #1
   \setlength{\leftmargin}{\cslhangindent}
   \setlength{\itemindent}{-1\cslhangindent}
  \fi
  % set entry spacing
  \setlength{\itemsep}{#2\baselineskip}}}
 {\end{list}}
\usepackage{calc}
\newcommand{\CSLBlock}[1]{\hfill\break\parbox[t]{\linewidth}{\strut\ignorespaces#1\strut}}
\newcommand{\CSLLeftMargin}[1]{\parbox[t]{\csllabelwidth}{\strut#1\strut}}
\newcommand{\CSLRightInline}[1]{\parbox[t]{\linewidth - \csllabelwidth}{\strut#1\strut}}
\newcommand{\CSLIndent}[1]{\hspace{\cslhangindent}#1}



\setlength{\emergencystretch}{3em} % prevent overfull lines

\providecommand{\tightlist}{%
  \setlength{\itemsep}{0pt}\setlength{\parskip}{0pt}}



 


\KOMAoption{captions}{tableheading}
\makeatletter
\@ifpackageloaded{caption}{}{\usepackage{caption}}
\AtBeginDocument{%
\ifdefined\contentsname
  \renewcommand*\contentsname{Table of contents}
\else
  \newcommand\contentsname{Table of contents}
\fi
\ifdefined\listfigurename
  \renewcommand*\listfigurename{List of Figures}
\else
  \newcommand\listfigurename{List of Figures}
\fi
\ifdefined\listtablename
  \renewcommand*\listtablename{List of Tables}
\else
  \newcommand\listtablename{List of Tables}
\fi
\ifdefined\figurename
  \renewcommand*\figurename{Figure}
\else
  \newcommand\figurename{Figure}
\fi
\ifdefined\tablename
  \renewcommand*\tablename{Table}
\else
  \newcommand\tablename{Table}
\fi
}
\@ifpackageloaded{float}{}{\usepackage{float}}
\floatstyle{ruled}
\@ifundefined{c@chapter}{\newfloat{codelisting}{h}{lop}}{\newfloat{codelisting}{h}{lop}[chapter]}
\floatname{codelisting}{Listing}
\newcommand*\listoflistings{\listof{codelisting}{List of Listings}}
\makeatother
\makeatletter
\makeatother
\makeatletter
\@ifpackageloaded{caption}{}{\usepackage{caption}}
\@ifpackageloaded{subcaption}{}{\usepackage{subcaption}}
\makeatother
\usepackage{bookmark}
\IfFileExists{xurl.sty}{\usepackage{xurl}}{} % add URL line breaks if available
\urlstyle{same}
\hypersetup{
  pdftitle={Probabilistic Models of Media (Mis)Trust},
  colorlinks=true,
  linkcolor={blue},
  filecolor={Maroon},
  citecolor={Blue},
  urlcolor={Blue},
  pdfcreator={LaTeX via pandoc}}


\title{Probabilistic Models of Media (Mis)Trust}
\usepackage{etoolbox}
\makeatletter
\providecommand{\subtitle}[1]{% add subtitle to \maketitle
  \apptocmd{\@title}{\par {\large #1 \par}}{}{}
}
\makeatother
\subtitle{Final Project, DSAN 5650: Causal Inference for Computational
Social Science}
\author{Probabilistic Modeler}
\date{2026-08-10}
\begin{document}
\maketitle
\begin{abstract}
Testing 123
\end{abstract}


\subsection{Introduction}\label{introduction}

I shall deal only briefly with an important variant of the seesaw coup.
This is a coup made by members of a group formerly well represented in
the state apparatus in order to retrieve their advantageous but
declining position.

Rare is the regime in Asia or Africa that is able to accord position and
influence to various ethnic groups in proportions generally acceptable
across group lines. Most governments are ethnically unbalanced. They
slight some groups; they lean toward others. Many such governments,
moreover, ride to power on promises to curtail the influence or
privilege of particular ethnic groups.

A group so excluded from its former position is unlikely to accept the
exclusion equanimously. Where such a group is heavily represented in the
officer corps, the regime is confronted with a difficult choice. To
purge these officers is to increase ethnic discontent, perhaps to bring
on a coup before the purging is complete, and to risk spreading
discontent among other officers who sense that careers abruptly ended
today for one reason can be equally abruptly ended tomorrow for other
reasons. On the other hand, to permit a potential focus of military
discontent to remain unchecked is also to risk a coup.

A coup to retrieve the lost or declining position of an ethnic group is
not inevitable. Such a coup depends on a particular configuration of
circumstances.

In the first place, groups that have recently been excluded from a
position of preeminence are more likely to be advanced ethnic groups:
disproportionately educated, urban, willing to migrate from their home
area, and stereotyped as diligent, industrious, energetic, and
intelligent. Typically, these groups are disproportionately represented
in educational institutions, in commerce and the professions, and in the
civil service. But such groups are not necessarily significantly
represented in the officer corps--unless educational prerequisites for
entry into military academies have been enforced and the military has
been a career affording salary and status at least roughly comparable to
civilian government service. Where these conditions are met,
representation of an excluded group in the officer corps may be
sufficient to make a retrieval coup possible.

Moreover, retrieval coups are mixed-motive coups, because a declining or
excluded group is unlikely to be able to bring off a coup all by itself.
For a retrieval coup to occur, officers who are not members of that
ethnic group must also be convinced that the exclusion threatens to
affect their own interests, or they must have their own reasons for
turning against the regime.

Two coups in Ghana illustrate these points. Both involved a complex mix
of motives, in which ethnicity played an important but not overriding
role.

As we have seen, the Ghanaian officer corps was drawn from the
well-educated Ewe, Ga, and Fanti groups. Nkrumah's electoral support was
also somewhat skewed. His Convention People's Party had been out-polled
in 1956 by the National Liberation Movement in Ashanti and by the
Northern People's Party in the North. The Ewe areas of the Trans Volta
were divided between the CPP and the Togoland Congress. No sooner had
independence been achieved than the regime was also confronted with an
Ewe separatist movement in the Trans Volta region and a militant Ga
movement in Accra.\footnote{See Smock and Smock (1975), 70.} Nkrumah
moved to suppress the dissidents, and the suppression increased after he
learned of conspiratorial approaches from opposition politicians to army
officers, principally of Ga and Ewe origin, in 1958. A shift in officer
recruitment away from Ga and Ewe was the result of these events.

Nkrumah offended his officers by tampering with their prerogatives. He
introduced party cells in the army, established a President's Own Guard
Regiment as a counterweight to the regular army, then favored the
Regiment in budgetary allocations, and played favorites in certain
conspicuous appointments and dismissals of officers.\footnote{First
  (1970), 197-200; Leys (1969), 235.}

Nkrumah's policy of using his emergency powers against opposition ethnic
leaders only increased ethnic grievances. The Ashanti did not forget
their humiliation at Nkrumah's hand. They complained of the domination
of Ashanti by the CPP's Southern leaders.\footnote{Welch (1970), 165.}
An advanced group from an unusually poor region, the Ewe were
particularly dependent on government largess for their survival. Bitter
at Nkrumah's punitive policy and the force with which he put down their
rebellion, the Ewe never accommodated themselves to CPP rule. Nor did
Nkrumah trust them; not a single Ewe was appointed to his cabinet after
1961. The three men who first plotted Nkrumah's overthrow in 1966 were
all Ewe. Of the eight members of the National Liberation Council (NLC)
established after the coup, three were Ewe and two were Ga.

The prominent role of Ewe in the coup and in the NLC gave rise to fears
of Ewe favoritism. During the NLC period, 1966---69, various Akan
groups---including some, such as the Ashanti and Fanti, that earlier had
been mutually suspicious---began to forge close links. What gave them a
new sense of cohesion was the alleged threat from the Ewe. Akans tended
to see the NLC as dominated by the Ewe and Ga.\footnote{Smock and Smock
  (1975), 203-205, 240-241.} When a countercoup failed in 1967, taking
the lives of three officers, all of them Ewe, the NLC felt obliged to
deny the ``wicked rumor'' that the move had been ``planned by Ashantis
and Fantis against Gas and Ewes.''\footnote{Quoted in Leys (1969), 243;
  First (1970), 403.}

The prospect of a return to civilian rule only deepened cleavages. Some
of the principal NLC members had ties to civilian politicians of the
same ethnic group. The leading NLC Ewe, John Harlley and A. K. Deku,
were close to the prominent Ewe leader K. Gbedemah. General J. A. Ankrah
and police commissioner John Nunoo, both Ga, were involved in efforts to
form a political party based on Ga support. Brigadier A. A. Afrifa, an
Ashanti, was closely linked to the Ashanti politician K. A. Busia. As
the end of military rule grew nearer, the Akan-Ewe rivalry, epitomized
by the antipathy between Afrifa and Harlley, was expressed in official
actions impinging on the political interests of the respective groups.
Ewe officers obtained, under threat of force, revocation of a decree
issued by Afrifa that would have banned from partisan activity a large
number of former CPP politicians, including Gbedemah.\footnote{First
  (1970), 405; Welch (1970), 200.} The removal of the ban on Gbedemah
fanned anti-Ewe sentiment. Shortly before the 1969 election, the NLC
removed two Ewe brigadiers from troop commands.

When civilian politics was restored in 1969, parties quickly
crystallized along ethnic lines. The rivalry of Harlley and Afrifa was
supplanted by the rivalry of their respective civilian ethnic kinsmen,
Gbedemah and Busia. Gbedemah's party appealed principally to the Ewe.
Nearly every other area voted for Busia's Progress Party, in large part
due to ``fears of Ewe domination of national
government\ldots{}''\footnote{Welch (1970), 220.}

Busia in office proved distinctly unfriendly to the Ewe. No Ewe served
in his cabinet, and legal action was even taken to unseat Gbedemah. The
list of 568 civil servants and policemen dismissed in 1970 contained
more than a fair share of Ewe names, and Ewe were pushed out of command
positions in the armed forces. By the end of 1971, only one Ewe was
still serving in a senior army position.\footnote{Smock and Smock
  (1975), 246-247; Austin and Luckham (1975), 304; Philippe Decraene,
  ``Ghana Coup,'' \emph{Le Monde} (English ed.), Jan 22, 1972.}

Once again, Ewe officers moved to retrieve their position. This they
did, as before, in conjunction with others acting (as no doubt they
themselves were) out of diverse motives. Two of the four organizers of
the 1972 coup were Ewe majors. Reduced military budgets, inequity in
promotions due to participation or nonparticipation in the coup against
Nkrumah, and civilian interference in army affairs---including punitive
action taken against Ewe officers---all appear to have had a part in the
decision to depose Busia's government. Most conspicuously
overrepresented in the new government's National Redemption Council
(NRC) were two identifiable groups: officers whose careers had stagnated
in the post-Nkrumah period and Ewe officers. At the outset, almost half
of the council was Ewe.\footnote{Austin and Luckham (1975), 308-309,
  312n18.}

A somewhat similar case could be made for the attempted coup in Sri
Lanka in 1962, which followed discriminatory action against Tamils and
Christians.\footnote{Horowitz (1980).} Both followed expansions of
political participation unaccompanied by drastic change in the
composition of the officer corps. Both regimes threatened the officer
corps on ethnic, social class, and organizational fronts simultaneously.
The concatenation of these threats linked officers who had otherwise
disparate interests, causing them to see the regimes' actions as part of
a coherent pattern of malevolent conduct. The essential ingredient of a
retrieval coup is simultaneous regime action that facilitates the
creation of the interethnic military coalition. For this purpose, the
combination of ethnic discrimination and tampering with military careers
is particularly likely to incite officers. Both can be subsumed under
the rubric of ``favoritism,'' which offends the conceptions of merit,
order, and regularity that are thought to underlie military recruitment,
tenure, and promotion.

Whereas a full seesaw coup is preceded by polarization of the polity
into two antagonistic ethnic clusters, which make mutually exclusive
claims on power, the retrieval coup aims merely to restore the former
position of a recently excluded group. The retrieval coup occurs in an
ethnically less explosive atmosphere, and the coupmakers are typically
multiethnic in composition. Ethnic retrieval is the goal for only some
of them. In the absence of polarized ethnic equilibrium, then, exclusion
of a single ethnic group need not split the army. In fact, it may
increase the cohesion of the army and make possible military
intervention based on a coalescence of differing motives.

Cohesion sufficient to make a coup, however, does not assure enough
cohesion to maintain interethnic alliances once the coup is made. After
both Ghanaian coups, reinclusion of the Ewe was aborted. Shortly after
the 1966 coup, Brigadier Afrifa propounded the historically dubious
proposition that the Ashanti and the Ewe were traditional
allies.\footnote{Afrifa (1966), 40.} As we have seen, it was not long
before Ewe and Ashanti members of the NLC fell out among themselves. By
1969, the Ewe were left to contest the elections practically alone in
their own party, and they were decisively defeated. A more serious split
followed the 1972 coup. Reflecting widespread distaste for the anti-Ewe
policy of the Busia regime, Colonel Acheampong said the coup was made
partly to combat ``the politics of tribalism.''\footnote{Quoted in Smock
  and Smock (1975), 249.} Again, however, the military council divided
along Akan-Ewe lines, Akan expressing discontent at what they felt was
inordinate Ewe influence in the NRC. This coincided with the reemergence
of Ewe separatism, which ultimately received some support from Togo. By
1975, the pressures on Acheampong to push the Ewe out were apparently
quite strong---some sources speak of an ultimatum from senior Akan
officers---and Acheampong purged the leading Ewe officers who had
conspired with him in 1972. Some months later, a number of Ewe former
officers and N.C.O.s were arrested and convicted of plotting a
coup.\footnote{For these events, see Legum, ed., \emph{Africa
  Contemporary Record, 1975-76}, p.~B694; \emph{West Africa}, Oct.~20,
  1975, p.~1257; \emph{ibid.}, Oct.~27, 1975, p.~1287; \emph{ibid.},
  Dec.~1, 1975, p.~1460; \emph{ibid.}, Jan.~5, 1976, p.~21;
  \emph{ibid.}, July 12, 1976, p.~981.} This was followed by another
unsuccessful Ewe plot in 1976. By this time, the Ewe were right back
where they were under the Busia regime: politically excluded despite
their significant representation in the officer corps.

In the three and a half years from July 1978 through December 1981,
Ghana experienced four more changes of government, three of them through
coups. In these kaleidoscopic events, the position of the Ewe was not
the principal issue, but the issue did not go away either. Jerry
Rawlings, who led coups in 1979 and 1981, is an Ewe. The first of his
regimes was disproportionately composed of Ga and Ewe, and it seemed to
engender resistance on these grounds among Akans. The elected government
that succeeded it embarked on a campaign against Rawlings and the
leading Ewe members of his council, ``fueling Ewe discontents to new
heights.''\footnote{Chazan (1982), 461-485, at 483.} After Rawlings'
second coup, in 1981, he cut back on the prominence of the Ewe, but, by
some accounts, Akan-Ewe conflicts nonetheless proceeded
apace.\footnote{\emph{ibid.}}

The Ewe thus tend to find themselves sharing power when the military
takes power and excluded soon after. Periods of civilian rule generally
leave the Ewe unrepresented, while periods of military rule serve to
sharpen political differences between the Ewe and the various Akan
groups. Ghanaian governments seem chronically unable to keep both Ewe
and Akans included on a stable basis.

\subsection{Conclusion}\label{conclusion}

\subsection{References}\label{references}

\protect\phantomsection\label{refs}
\begin{CSLReferences}{1}{0}
\bibitem[\citeproctext]{ref-afrifa_ghana_1966}
Afrifa, Akwasi Amankwaa. 1966.
\emph{\href{https://books.google.com?id=5IQhAAAAMAAJ}{The {Ghana Coup},
24th {February} 1966}}. Cass.

\bibitem[\citeproctext]{ref-austin_politicians_1975}
Austin, Dennis, and Robin Luckham. 1975.
\emph{\href{https://books.google.com?id=oa3rAgAAQBAJ}{Politicians and
{Soldiers} in {Ghana} 1966-1972}}. Routledge.

\bibitem[\citeproctext]{ref-chazan_ethnicity_1982}
Chazan, Naomi. 1982. {``\href{https://doi.org/10.2307/2149995}{Ethnicity
and {Politics} in {Ghana}}.''} \emph{Political Science Quarterly} 97
(3): 461--85.

\bibitem[\citeproctext]{ref-first_barrel_1970}
First, Ruth. 1970.
\emph{\href{https://books.google.com?id=8P0EAQAAIAAJ}{The {Barrel} of a
{Gun}: {Political Power} in {Africa} and the {Coup D}'état}}. Allen
Lane.

\bibitem[\citeproctext]{ref-horowitz_coup_1980}
Horowitz, Donald L. 1980.
\emph{\href{https://books.google.com?id=ZbP_AwAAQBAJ}{Coup {Theories}
and {Officers}' {Motives}: {Sri Lanka} in {Comparative Perspective}}}.
Princeton University Press.

\bibitem[\citeproctext]{ref-leys_politics_1969}
Leys, Colin. 1969.
\emph{\href{https://books.google.com?id=FkcJbkpXufYC}{Politics and
{Change} in {Developing Countries}: {Studies} in the {Theory} and
{Practice} of {Development}}}. Cambridge University Press.

\bibitem[\citeproctext]{ref-smock_politics_1975}
Smock, David R., and Audrey C. Smock. 1975.
\emph{\href{https://doi.org/10.2307/2064848}{The {Politics} of
{Pluralism}: {A Comparative Study} of {Lebanon} and {Ghana}.}}

\bibitem[\citeproctext]{ref-welch_soldier_1970}
Welch, Claude Emerson. 1970.
\emph{\href{https://books.google.com?id=xAC5AAAAIAAJ}{Soldier and
{State} in {Africa}: A {Comparative Analysis} of {Military Intervention}
and {Political Change}}}. Northwestern University Press.

\end{CSLReferences}




\end{document}
